
\epigraph{``\textit{Non-reproducible single occurrences are of no significance to
science.}''}{---\citeauthor{popper2005logic},\citeyear{popper2005logic} \cite{popper2005logic}}
Los experimentos realizados en esta tarea se llevaron a cabo en un entorno controlado para garantizar la reproducibilidad de los resultados. La infraestructura utilizada es la siguiente:

\begin{itemize}
    \item \textbf{Hardware:} Procesador	11th Gen Intel(R) Core(TM) i5-1135G7 @ 2.40GHz   2.42 GHz, 8 GB RAM, almacenamiento SSD Intel Optane.
    \item \textbf{Software:} Sistema operativo Ubuntu 22.04 LTS corriendo bajo WSL (Windows Subsystem for Linux) sobre Windows 10. Compilador \texttt{g++ 11.4.0} con soporte para \texttt{-std=c++17}. Scripts complementarios en Python 3.10.12.
\end{itemize}

Ambos algoritmos fueron compilados utilizando un \texttt{makefile}, lo que permite replicar el proceso sin ambigüedades. Las mediciones se realizaron ejecutando ambos programas sobre un conjunto de datos fijo, generado por un script controlado, que produce un archivo con 10 casos de prueba variados en tamaño, estructura y contenido. Cada ejecución procesa exactamente los mismos datos de entrada, asegurando una comparación justa entre los paradigmas.

Se midió el tiempo de ejecución de cada caso utilizando la librería estándar \texttt{<chrono>} en C++, y se estimó el uso de memoria con base en las estructuras activas utilizadas por cada algoritmo. Los resultados se almacenaron en archivos de texto para su posterior graficación.

El propósito de los experimentos es comparar de forma empírica la eficiencia y escalabilidad de dos estrategias algorítmicas para resolver el mismo problema, en condiciones equivalentes y sobre entradas diseñadas para mostrar diferencias significativas de rendimiento.
